\subsection{HPX -- High Performance ParallelX}
\label{sec:frameworks:hpx}

\begin{quotebox}{\href{https://hpx-docs.stellar-group.org/latest/html/index.html\#what-is-hpx}{HPX Documentation}}
    HPX is a C++ Standard Library for Concurrency and Parallelism. It implements all of the corresponding facilities as defined by the C++ Standard. Additionally, in HPX we implement functionalities proposed as part of the ongoing C++ standardization process. We also extend the C++ Standard APIs to the distributed case. HPX is developed by the STE||AR group.
\end{quotebox}

\acrfull{hpx}~\cite{kaiser_hpx_2020} is high-performance many-task asynchronous tasking library. 
Its development started in 2007~\cite{the_stear_group_history_nodate} with the current version being v1.10.0 from March 2024~\cite{kaiser_stellar-grouphpx_2024}.
Although \gls{hpx} is not standardized where possible, its \gls{api} is closely modeled after the \cpp standard \gls{api}. 

The execution model of \gls{hpx} is based on a task graph that is implicitly built in the source code and later executed asynchronously. 
It implements custom schedulers that also incorporate work-stealing to automatically improve the load-balancing of applications using \gls{hpx}. 
This is especially useful for non-uniform problems that rely on good load balancing. 
\gls{hpx} also supports distributed memory architectures using an \gls{agas} and different accelerators using executors. 
Note that \gls{hpx} only portably incorporates kernel launches in its internal task graph; the actual compute kernels must still be written using other frameworks like \gls{cuda}, \gls{hip}, Kokkos, or SYCL.  
Additionally, \gls{hpx} supports the \cpp{17} standard library parallel algorithms (also see~\autoref{sec:frameworks:stdpar}. 
In the remainder of this thesis, only these standard algorithms are used with \gls{hpx}. 

\todo{mention that parallel algorithms were only implemented after the \cpp draft}

\begin{listing}
    \begin{minted}{cpp}
std::vector<int> indices(N);
std::iota(indices.begin(), indices.end(), 0);

hpx::for_each(hpx::execution::par_unseq, indices.begin(), indices.end(), [&](const int idx) {
    Y[idx] = alpha * X[idx] + Y[idx];
});
    \end{minted}
    \caption{%
    A simple DAXPY example using \gls{hpx} parallel loops.}
    \label{code:hpx_example}
\end{listing}

An example of these \cpp{17} standard parallel algorithms can be seen in \autoref{code:hpx_example}. 
The only difference compared to \autoref{code:stdpar_example} is the usage of the \code{hpx::} namespace qualifiers instead of the \code{std::} ones. 
In this code, the \code{for_each} loop is automatically parallelized across all available \gls{cpu} cores. 


%%% ADVANTAGES AND DISADVANTAGES
\AdvantagesAndDisadvantages{hpx}{
    \item supports algorithms close to \cpp standard
    \item task-based abstraction more suitable for some problems
}{
    \item not standardized
    \item can only launch kernels, other frameworks still \textbf{necessary} to write the actual compute kernels 
}
